\newpage
\section{Turbo Codes}

\subsection{Generating the Trellis}
The sequence of output bits is generated by computing the parity of the 
product between the generator polynomial and the contents of the shift register. 
If the shift register has length $k$, the encoder requires $2^{k-1}$ distinct states, 
since each new bit is obtained by appending the input bit to the register and 
then applying the parity operation. For every state and input pair, we compute both 
the next state of the register and the corresponding output parity bit. 
Using this information, we construct the trellis diagram over the full sequence 
of input bits.\\
Below is an example trellis with a shift register of length 3, generators 5 and 7, and a sequence of 4 bits. (Start state for any sequence is 00)

\begin{center}
\begin{tikzpicture}[
    ->, >=Stealth, auto, node distance=3.5cm,
    state/.style={circle, draw, minimum size=0.7cm, font=\scriptsize, inner sep=1pt},
    edge label/.style={font=\tiny, midway, above, sloped, inner sep=1pt},
    every edge/.style={draw, thick}
]

% Time step labels
\node[font=\scriptsize] at (-0.5, 4.5) {$t=0$};
\node[font=\scriptsize] at (3, 4.5) {$t=1$};
\node[font=\scriptsize] at (6.5, 4.5) {$t=2$};
\node[font=\scriptsize] at (10, 4.5) {$t=3$};

% States at t=0
\node[state] (S0_0) at (0, 4) {$00$};
\node[state] (S1_0) at (0, 3) {$01$};
\node[state] (S2_0) at (0, 2) {$10$};
\node[state] (S3_0) at (0, 1) {$11$};

% States at t=1
\node[state] (S0_1) at (3.5, 4) {$00$};
\node[state] (S1_1) at (3.5, 3) {$01$};
\node[state] (S2_1) at (3.5, 2) {$10$};
\node[state] (S3_1) at (3.5, 1) {$11$};

% States at t=2
\node[state] (S0_2) at (7, 4) {$00$};
\node[state] (S1_2) at (7, 3) {$01$};
\node[state] (S2_2) at (7, 2) {$10$};
\node[state] (S3_2) at (7, 1) {$11$};

% States at t=3
\node[state] (S0_3) at (10.5, 4) {$00$};
\node[state] (S1_3) at (10.5, 3) {$01$};
\node[state] (S2_3) at (10.5, 2) {$10$};
\node[state] (S3_3) at (10.5, 1) {$11$};

% Transitions from t=0 to t=1
\path (S0_0) edge[edge label, above] node {$0/0$} (S0_1);
\path (S0_0) edge[edge label, bend left=5, above] node {$1/1$} (S1_1);
\path (S1_0) edge[edge label, bend right=25, above] node {$0/1$} (S3_1);
\path (S1_0) edge[edge label, bend left=15, above] node {$1/0$} (S2_1);
\path (S2_0) edge[edge label, above] node {$0/0$} (S1_1);
\path (S2_0) edge[edge label, bend left=15, above] node {$1/1$} (S0_1);
\path (S3_0) edge[edge label, bend left=20, above] node {$0/1$} (S2_1);
\path (S3_0) edge[edge label, above] node {$1/0$} (S3_1);

% Transitions from t=1 to t=2
\path (S0_1) edge[edge label, above] node {$0/0$} (S0_2);
\path (S0_1) edge[edge label, bend left=5, above] node {$1/1$} (S1_2);
\path (S1_1) edge[edge label, bend right=25, above] node {$0/1$} (S3_2);
\path (S1_1) edge[edge label, bend left=15, above] node {$1/0$} (S2_2);
\path (S2_1) edge[edge label, above] node {$0/0$} (S1_2);
\path (S2_1) edge[edge label, bend left=15, above] node {$1/1$} (S0_2);
\path (S3_1) edge[edge label, bend left=20, above] node {$0/1$} (S2_2);
\path (S3_1) edge[edge label, above] node {$1/0$} (S3_2);

% Transitions from t=2 to t=3
\path (S0_2) edge[edge label, above] node {$0/0$} (S0_3);
\path (S0_2) edge[edge label, bend left=5, above] node {$1/1$} (S1_3);
\path (S1_2) edge[edge label, bend right=25, above] node {$0/1$} (S3_3);
\path (S1_2) edge[edge label, bend left=15, above] node {$1/0$} (S2_3);
\path (S2_2) edge[edge label, above] node {$0/0$} (S1_3);
\path (S2_2) edge[edge label, bend left=15, above] node {$1/1$} (S0_3);
\path (S3_2) edge[edge label, bend left=20, above] node {$0/1$} (S2_3);
\path (S3_2) edge[edge label, above] node {$1/0$} (S3_3);

% Legend
\node[font=\scriptsize, align=left, draw, inner sep=3pt] at (10.5, 0) {
    $S_t$: State at time \\ 
    (last 2 shift register bits) \\
    Edge: input/parity \\
    Generators: $G_1=101$, $G_2=111$
};

\end{tikzpicture}
\end{center}

\subsection{Viterbi}
The algorithm begins in state $0$ with an initial cost of $0$, while all other 
states are assigned an infinite cost. Here, the cost represents the negative 
log likelihood of the most probable sequence leading to a particular state. 
\\During the forward pass, the algorithm iterates over each time step 
(corresponding to one bit of the original sequence) and evaluates all possible 
states at that step. For every possible input bit, it computes the likelihood of 
observing both the systematic and parity bits. These likelihoods are converted 
to costs and added to the cumulative cost of reaching the current state. If this 
new path yields a lower cost than any previously recorded path to the next state, 
the cost is updated along with a backpointer to record the optimal predecessor. 
\\In the backtracking phase, the algorithm identifies the final state with the 
minimum cumulative cost (equivalently, the maximum likelihood). Starting from 
this state, it follows the backpointers in reverse through the trellis, thereby 
reconstructing the most probable sequence of input bits.

The time complexity of the algorithm is $O(T * S)$ where $T$ is the length of the original string or the number of timestamps and $S$ is the total number of states (in this case $2^{k - 1}$) where k is the length of the shift register.

\subsection{BCJL and Turbo Codes}
In turbo decoding, a posteriori probabilities (APPs) are estimated with the BCJR algorithm, starting from uniform priors ($P(b=0)=P(b=1)=\tfrac{1}{2}$). For each received sequence, BCJR computes log-likelihood ratios (LLRs):
\[
    LLR_i = \log \frac{P(u_i=1 \mid \mathbf{y})}{P(u_i=0 \mid \mathbf{y})},
\]
where $u_i$ is the transmitted bit and $\mathbf{y}$ the channel output.
\\Each decoder then forms \emph{extrinsic information} by removing the prior 
contribution:
\[
    LLR^{\text{extrinsic}}_i = LLR_i - LLR^{\text{prior}}_i,
\] 
\[
    LLR^{\text{prior}}_i = \log \frac{P(u_i=1)}{P(u_i=0)}.
\]
These extrinsic values are interleaved and passed to the other decoder as new priors. Iterating this exchange gradually refines the APPs.
\\Convergence is monitored via the maximum LLR change,
\[
    \Delta = \max_i \left| LLR^{(t)}_i - LLR^{(t-1)}_i \right|,
\]
stopping early if $\Delta$ falls below a threshold, or otherwise after a fixed maximum number of iterations.

\subsection{Results}
% Run the code multiple times, for sizes 1000, 2000, 5000, 10000
% Write down the average BER (bit error rate)
We compared the decoding performance of three algorithms: the Viterbi decoder, 
the BCJR (MAP) decoder, and the Turbo decoder. The cumulative average BERs across 
all experiments are summarised below (and in results.txt):

\begin{itemize}
    \item Viterbi decoder: \textbf{3.90\%}
    \item BCJR decoder: \textbf{3.61\%}
    \item Turbo decoder: \textbf{0.0067\%}
\end{itemize}
In terms of BER, BCJL performs slightly better than Viterbi on average. This might be because BCJL gives the probability for each bit individually. The average length of a matched string is better in Viterbi, probably because Viterbi outputs the most probable string.
\\Since the turbo code always has almost zero errors, applying Viterbi on the final APP helps but does not improve the BER to any significant degree.